\section{Umsetzung der Integration mit Power Query}\label{sec:integration}

% Zielumfang: ca. 5 Seiten – der Kern der Arbeit. Jede Entscheidung begründen, Alternative nennen. Screenshots sparsam, M-Code-Auszüge dort, wo der Klickpfad nicht reicht.

\subsection{Aufbau der Abfragen}\label{sec:aufbau}

\todo{Abbildung: Abfrageabhängigkeiten aus dem Power-Query-Editor (Ansicht \enquote{Abfrageabhängigkeiten}).}

Die Mappe ist in Schichten gegliedert, die der Staging-Idee des Skripts folgen \skript{8.2}:
\begin{description}
  \item[Rohdaten] \pq{KUNDE}, \pq{customers}, \pq{Accounts}, \dots\ -- je Quelldatei eine Abfrage, nur Import mit korrekter Codepage, Trennzeichen und Gebietsschema, keine Logik.
  \item[Hilfsfunktionen und Mappings] \pq{fnClean}, \pq{fnNormEmail}, \pq{fnNormStrasse}, \pq{fnNormName}; Mapping-Tabellen \pq{mapLand}, \pq{mapKundentyp}, \pq{mapRegion}, \pq{mapMwSt}; Parameter \pq{pSchwelle}, \pq{pKursCHF}.
  \item[Harmonisierung (\_h)] \pq{KUNDE\_h}, \pq{SHOP\_h}, \pq{CRM\_h} -- je Quelle ein einheitliches Zwischenschema mit \pq{LandISO}, \pq{PLZ\_norm}, \pq{Kundentyp}, \pq{Matchkey}, \pq{Name\_norm}, \pq{Email\_norm}.
  \item[Kundenintegration (\_k)] \pq{KUNDE\_k} (Basis, vergibt \pq{KundeKey}), \pq{ERP\_Dubletten}, \pq{CRM\_k}, \pq{SHOP\_k}, \pq{MAP\_SCHLUESSEL}.
  \item[Ziel] \pq{DIM\_KUNDE}, \pq{DIM\_KUNDE\_ADRESSE}, \pq{DIM\_VERKAEUFER}, \pq{DIM\_PRODUKT}, \pq{FAKT\_VERKAUF}.
  \item[Prüfung] \pq{CHECK\_Sollwerte}, \pq{CHECK\_Detail}, \pq{CHECK\_Fakten}.
\end{description}

\todo{Integrationsstrategie begründen (Checkliste 35): leiterartig ERP $\rightarrow$ +CRM $\rightarrow$ +Shop, weil das ERP die führenden Stammdaten und die Verkäufer hat; Alternative n-är nennen.}

\subsection{Syntaktische Harmonisierung}\label{sec:syntax}

\todo{Checkliste 1--4, 38, 39, 45: Codepage 1252 vs. 65001; Typ ändern mit Gebietsschema de-DE/en-US; drei Datumsformate; Zeitzone; Csv.Document mit QuoteStyle.Csv wegen eingebetteter Zeilenumbrüche (Nachweis: 124 physische Zeilen $\rightarrow$ 90 Datensätze); fnClean gegen nachgestellte/doppelte/geschützte Leerzeichen VOR dem Matchkey; Mojibake-Rückübersetzung.}

\begin{lstlisting}[language=PowerQueryM, caption={Hilfsfunktion \pq{fnClean}: unsichtbare Zeichen und Mehrfach-Leerzeichen entfernen}, label={lst:fnclean}]
(s as nullable text) as nullable text =>
let
    t0 = if s = null then null else Text.Clean(s),
    t1 = if t0 = null then null else Text.Replace(t0, Character.FromNumber(160), " "),
    t2 = if t1 = null then null else Text.Remove(t1, {Character.FromNumber(173), Character.FromNumber(8203)}),
    t3 = if t2 = null then null else Text.Combine(List.Select(Text.Split(t2, " "), each _ <> ""), " "),
    t4 = if t3 = "" then null else t3
in
    t4
\end{lstlisting}

\subsection{Kodierungen, Schema und Datenfehler}\label{sec:schema}

\todo{Checkliste 6--12, 15--20, 23--26, 41--44, 46, 48: Ja/Nein-Kodierungen, Anrede, NULL-Kodierungen (-1, 31.12.9999), Länder-Mapping (stark äquivalent) vs. Kundentyp-Mapping (schwach äquivalent, kleinster gemeinsamer Nenner), is\_club $\rightarrow$ Kundentyp, Lieferart Entität $\rightarrow$ Attribut, Synonyme/Homonyme umbenennen, Titel aus Name1, Straße/Hausnummer teilen und vertauschte Felder zurücktauschen, PLZ mit führender Null, Land aus PLZ ableiten, MM/YYYY mit Flag\_Datum\_ungenau, Telefon E.164.}

\todo{Ein bis zwei Screenshots: Mapping-Tabelle mapKundentyp mit Merge; bedingte Spalte Strasse\_ok/Hausnr\_ok.}

\subsection{Kundenintegration: Record Linkage als Kaskade}\label{sec:linkage}

\todo{Kernstück. Checkliste 21, 22, 27--31, 47.}

Die drei Quellen haben keinen gemeinsamen Schlüssel: ERP und CRM teilen \pq{ERP\_KdNr} nur bei 60\,\% der Accounts und teils fehlerhaft, Shop und CRM teilen die E-Mail, das ERP führt keine E-Mail. Deshalb wird je Quelle eine Kaskade durchlaufen, die den sichersten verfügbaren Abgleich zuerst versucht und die Matchmethode je Zeile protokolliert:

\begin{table}[H]
  \centering\small
  \caption{Match-Kaskade je Quelle (Stand Messung 3b)}
  \label{tab:kaskade}
  \begin{tabular}{@{}c>{\raggedright\arraybackslash}p{4.3cm}>{\raggedright\arraybackslash}p{4.6cm}>{\raggedright\arraybackslash}p{4.0cm}@{}}
    \toprule
    Stufe & Kriterium & CRM $\rightarrow$ ERP & Shop $\rightarrow$ ERP/CRM \\
    \midrule
    1 & ERP-Schlüssel + Namensprüfung & 38 Treffer, 1 Tippfehler abgewiesen & 24 Treffer \\
    2 & E-Mail normalisiert & entfällt (ERP ohne E-Mail) & 51 Treffer (CRM) \\
    3 & Adresse exakt über \pq{Matchkey} & 18 Treffer & 0 (bereits über 1/2 gefunden) \\
    4 & Fuzzy auf \pq{Name\_norm} + Nachregel & 5 Treffer & 0 \\
    -- & neu (eigener \pq{KundeKey} ab 1001 / 2001) & 49 (15 Leads, 30 nur Shop+CRM, 4 Umzieher) & 65 (reine Privatkunden) \\
    \bottomrule
  \end{tabular}
\end{table}

\paragraph{Matchkey.} \todo{Land + PLZ\_norm + fnNormStrasse + Hausnummer -- aus PA1 übernommen und begründen, warum Name nicht in den Schlüssel gehört.}

\paragraph{Schlüsselprüfung.} \todo{Warum der Fremdschlüssel nicht blind übernommen wird: Namensprüfung fängt den eingebauten Tippfehler (ERP\_KdNr 10368 zeigt auf einen anderen Kunden). Flag\_Schluessel\_Name.}

\paragraph{Fuzzy-Abgleich und Nachregel.} \todo{Table.FuzzyNestedJoin liefert keinen Ähnlichkeitswert und vervielfacht Zeilen (122 statt 115) -- je Zeile ein Treffer. Die fachliche Nachregel entscheidet: Verein $\rightarrow$ nur identischer Name\_norm (Ort steckt im Namen); Händler/Gastro $\rightarrow$ gleiche PLZ oder identischer Name bei gleichem Ort. Begründen mit den Beobachtungen aus der Messreihe.}

\begin{lstlisting}[language=PowerQueryM, caption={Auswahl des Fuzzy-Treffers über die fachliche Nachregel (Auszug aus \pq{CRM\_k})}, label={lst:fuzzy}]
FE = Table.AddColumn(FJ, "Key_Fuzzy", each
  let k = Table.SelectRows([FZ], (r) =>
        if [Kundentyp] = "Verein" then r[Name_norm] = [Name_norm]
        else (r[PLZ_norm] = [PLZ_norm]
              or (r[Name_norm] = [Name_norm]
                  and Text.StartsWith(Text.Lower([Ort]), Text.Lower(r[Ort])))))
  in if Table.IsEmpty(k) then null else k{0}[KundeKey])
\end{lstlisting}

\paragraph{Dubletten innerhalb einer Quelle.} \todo{ERP: Gruppieren nach Matchkey + Name\_norm, ältester Datensatz führend, beide KdNr auf einen KundeKey (MAP\_SCHLUESSEL). CRM: Dubletten teilen sich beim Vergeben neuer Keys einen Key. Befund Anlagedatum als Text $\rightarrow$ falscher führender Datensatz.}

\subsection{Erfolgsmessung gegen die Sollwerte}\label{sec:erfolgsmessung}

Die Integration wird nicht nur plausibilisiert, sondern gemessen: \pq{CHECK\_Sollwerte} vergleicht die vergebenen \pq{KundeKey}s mit der Ground Truth aus der Datengenerierung. Zwei Fehlerarten werden gezählt: \emph{zerrissen} (ein echter Kunde hat mehrere Keys) und \emph{verschmolzen} (ein Key enthält mehrere echte Kunden). \Cref{tab:messreihe} zeigt die Messreihe; das vollständige Protokoll steht in \cref{app:messprotokoll}.

\begin{table}[H]
  \centering\small
  \caption{Messreihe der Kundenintegration (200 eindeutige Kunden, \pq{pSchwelle} = 0,8 sofern nicht anders angegeben)}
  \label{tab:messreihe}
  \begin{tabular}{@{}c>{\raggedright\arraybackslash}p{6.3cm}rrrc@{}}
    \toprule
    Nr. & Stand & zerr. & verschm. & o. Map. & Fuzzy CRM/Shop \\
    \midrule
    1  & Kaskade Schlüssel(+Name) $\rightarrow$ E-Mail $\rightarrow$ Adresse $\rightarrow$ Fuzzy (Name+Land, PLZ-Regel außer Vereine) & 5 & 12 & 0 & 19 / 0 \\
    2  & wie 1, Quelldaten korrigiert: Namen je Ort eindeutig & 7 & 8 & 0 & 13 / 0 \\
    3a & wie 2, Fuzzy nur bei gleicher PLZ oder (Verein und identischer Name\_norm) & 7 & 3 & 0 & 5 / 0 \\
    3b & Verein $\rightarrow$ nur identischer Name\_norm; sonst gleiche PLZ oder (identischer Name und gleicher Ort) & \textbf{4} & \textbf{0} & \textbf{0} & 5 / 0 \\
    4  & wie 3b, Schwellwert 0,9 & 4 & 0 & 0 & 5 / 0 \\
    5  & wie 3b, Schwellwert 0,7 & 4 & 0 & 0 & 5 / 0 \\
    \bottomrule
  \end{tabular}
\end{table}

\todo{Erkenntnisse aus dem Messprotokoll in 3--4 Absätze: (1) Verschmelzungen durch gleichnamige Vereine in verschiedenen Städten -- Fuzzy gewichtet alle Zeichen gleich, Ortsanteil ist klein; (2) Precision/Recall-Zielkonflikt: die PLZ-Regel, die Verschmelzungen verhindert, blockiert Umzüge; bewusst zugunsten Precision entschieden, weil falsche Verschmelzungen in der Umsatzanalyse schwerer wiegen; (3) Schwellwert ist nach 3b wirkungslos -- der Fuzzy-Join ist nur Kandidatenlieferant, die fachliche Nachregel entscheidet; Folgerung: Ähnlichkeitsmaß ohne Nachregel ist für Firmennamen mit Ortsbestandteil ungeeignet; (4) Gegenbeispiel SHOP\_k: mit einem verlässlichen gemeinsamen Attribut (E-Mail) kamen Adresse und Fuzzy nie zum Zug.}

\todo{Die 4 verbleibenden Zerreißungen sind Umzieher ohne Schlüssel und E-Mail $\rightarrow$ Prüfliste für den Fachbereich (Mängelklasse 2 nach Tab.~8.1 im Skript). Alternative \enquote{identischer Name genügt} nennen und verwerfen.}

\subsection{Faktenseite und abgeleitete Kennzahlen}\label{sec:fakten}

\todo{Checkliste 13, 14, 26, 32--34, 37: komposite Schlüssel verketten, Produktgruppe aus sprechendem Schlüssel, ArtNr$\leftrightarrow$SKU, CHF$\rightarrow$EUR, Brutto$\rightarrow$Netto je Land über mapMwSt, Set$\rightarrow$Stück, Filter (gesperrt, cancelled/refunded) mit Zählung, Lieferdauer und Zahlungsverzug, Saison, Region, Zuordnung des Auftrags zur gültigen Adresse per Datumsbereich. Dublettenschutz: Shop-Aufträge nur aus einer Quelle (ShopRef).}

\todo{CHECK\_Fakten: Summe Umsatz je Quelle vor/nach Integration, Anzahl Positionen, Anteil ohne Verkäufer.}

\subsection{Prüfliste und Grenze zwischen automatisch und manuell}\label{sec:pruefliste}

\todo{Checkliste 36: Jede Checklistenzeile einer Mängelklasse zuordnen (automatisch / halbautomatisch / manuell, Tab.~8.1 im Skript). Prüfliste mit Zählung: unscharfe Matches ohne Nachregel-Treffer, Umzieher ohne Schlüssel, PLZ-Land-Konflikte, ungültige PLZ/E-Mail, vertauschte Felder mit Zusatz. Das ist der Nachweis, dass die Grenze bewusst gezogen wurde.}

\todo{Abschließend ein Pivot-Beispiel je Auswertungsfrage oder mindestens F1 und F5 als Screenshot -- \enquote{Verdeutlichen Sie den Integrationsablauf mit Beispieldaten}.}
